% ===================================================================
%  圖表第 12 號 — 火車站。--- 鐵路。--- 船。
%
%  Traduction, faite en 2026, en cantonais ecrit d'aujourd'hui, en
%  caracteres traditionnels, sur l'ido de texto/io. Regles de la
%  colonne : voir 10-toubiu-01.tex. Une seule numerotation.
%
%  LE BILLET DONNE ICI LE MOT LE PLUS CANTONAIS DU LIVRET, ET IL TIENT
%  EN UN SEUL CARACTERE :
%
%    飛 (9)     le billet    Pekin : 票
%
%  C'est l'anglais « fare », entre par Hong Kong comme 芝士 et
%  車厘子, mais reduit a UNE syllabe et devenu si ordinaire qu'il se
%  conjugue : 買飛 prendre un billet, 車飛 le billet de train, 飛價
%  le prix des places. Le caractere 飛 veut dire « voler » et n'a
%  aucun rapport : il ne sert ici qu'a porter le son.
%
%  LE RESTE DU TRAIN SUIT :
%
%    車卡 (37)   la voiture    Pekin : 车厢
%    路軌 (64)   les rails     Pekin : 铁轨
%    火車頭 (54) la locomotive Pekin : 机车
%    侍應 (35)   le garcon     Pekin : 服务员
%    月台 (53)   le quai
%
%  « 車卡 » vaut la peine d'etre note : 卡 est lui-meme un emprunt
%  (« car »), et le compose met donc un mot anglais derriere un mot
%  chinois. Le nord a preferé 车厢, « la caisse du char ».
%
%  DEUX BLOCS ONT DEMANDE L'INVERSION : l'employe et sa bascule
%  (t12-05-1), ou l'ido nomme l'homme avant la machine dont il a la
%  charge ; et le renvoi (4) du bureau du telegraphe (t12-02-2), qui
%  vient APRES le controleur a qui on le demande. Ce dernier avait
%  deja fait trebucher la colonne fr-CA sur la meme planche.
%
%  L'APPEL DE NOTE DU BLOC 19 EST EN DEMI-CHASSE — « (*) » et non
%  « （*） » — parce que html.py apparie sur ce caractere-la.
% ===================================================================

%%K t12-apar-1 apar
\VUsaut{4.0mm}
\VUtitre{15.0pt}{60}{圖表第 12 號}
\VUsaut{2.4mm}
\VUpk{11.4pt}{40}{火車站。--- 鐵路。--- 船。}
\VUsaut{3.4mm}
\textsc{註。} --- \textit{各國用嘅時間表都唔同，所以我哋攞}
\VUgras{法國表}\textit{嘅鐘數做例。}

%%K t12-01-1 p
1. --- 我啱啱去咗德國一趟。出發之前，我買咗本\VUgras{火車時間表}
\textsuperscript{(15)}，睇下啲車幾點開。查到最就手嗰班係夜晚九點喺巴黎開，
凌晨一點十三分到史特拉斯堡，我就執定行李；第二日，我就叫人送我去火車站。

%%K t12-02-1 p
2. --- 我行入個大候車大堂嗰陣，前面有位鄉下嚟嘅老\VUgras{神父}
\textsuperscript{(2)}，佢手上攞住個\VUgras{旅行袋} \textsuperscript{(3)}；佢後面
已經到咗好多\VUgras{旅客} \textsuperscript{(1)}。

%%K t12-02-2 p
我想打份\VUgras{電報} \textsuperscript{(5)}，就叫個\VUgras{查票員} \textsuperscript{(6)}
指個\VUgras{電報房} \textsuperscript{(4)} 畀我。

%%K t12-03-1 p
3. --- 入\VUgras{候車室} \textsuperscript{(7)} 之前，我去買咗張來回\VUgras{飛}
\textsuperscript{(9)}。

%%K t12-04-1 p
4. --- 我喺個\VUgras{售票窗} \textsuperscript{(8)} 冇等好耐，因為冇乜人。有個放假
出發嘅\VUgras{兵} \textsuperscript{(10)} 好人噉讓咗個位畀我，我就仲夠時間去問個
\VUgras{站長} \textsuperscript{(13)} 幾樣嘢——佢當時同負責巡查嘅\VUgras{督察}
\textsuperscript{(12)}、同埋當值嘅\VUgras{憲兵} \textsuperscript{(11)} 傾緊偈。

%%K t12-tit-1 sub
\VUpk{10.2pt}{40}{寄行李。}
\VUsaut{3.0mm}

%%K t12-05-1 p
5. --- 我喺個\VUgras{書報檔} \textsuperscript{(14)} 買咗份報紙，跟住去磅我啲
\VUgras{行李} \textsuperscript{(17)}——啲行李係個\VUgras{搬運佬} \textsuperscript{(16)} 啱啱
用架\VUgras{行李車} \textsuperscript{(19)} 推埋嚟嘅。個\VUgras{職員} \textsuperscript{(22)}——
負責架\VUgras{磅} \textsuperscript{(21)} 嗰個——話我知我個櫳三十五公斤，即係
\VUgras{超重}五公斤；超重就要喺行李\VUgras{窗口} \textsuperscript{(23)} 補錢。

%%K t12-06-1 p
6. --- 我早咗好多，得閒就睇下站入面啲旅客點行點走。有啲喺個
\VUgras{寄存處} \textsuperscript{(25)} 寄或者攞返啲唔重要嘅行李；有啲就睇緊個
\VUgras{時間表} \textsuperscript{(26)}。啱到站嘅旅客趕住去\VUgras{出口}
\textsuperscript{(32)}，手上攞住個\VUgras{皮箱} \textsuperscript{(24)}。有幾個喺\VUgras{領行李處} \textsuperscript{(27)} 等，
遞上張\VUgras{收據} \textsuperscript{(29)} 攞返自己啲櫳、\VUgras{木箱} \textsuperscript{(28)}
或者\VUgras{籮} \textsuperscript{(30)}。

%%K t12-07-1 p
7. --- \VUgras{啲稅局人員} \textsuperscript{(33)} 就仔細噉查啲包裹，睇下有冇嘢
要報。

%%K t12-08-1 p
8. --- 有啲旅客入咗個\VUgras{餐室} \textsuperscript{(34)}，個\VUgras{侍應}
\textsuperscript{(35)} 落足心機服侍佢哋。佢哋要快手啲，因為架\VUgras{火車}
\textsuperscript{(36)} 係特快，唔會喺站度停耐。

%%K t12-09-1 p
9. --- 之後我行去開車嗰邊月台，搵架直通嘅\VUgras{車卡} \textsuperscript{(37)}，
噉就唔使中途轉車。我上咗架有走廊嘅車，入面有煙民用嘅\VUgras{車廂}
\textsuperscript{(38)}，仲有一個係留畀「單身女士」嘅。

%%K t12-tit-2 sub
\VUpk{10.2pt}{40}{夜晚開車。}
\VUsaut{3.0mm}

%%K t12-10-1 p
10. --- 我踏上\VUgras{踏板} \textsuperscript{(40)}、閂埋道\VUgras{門}
\textsuperscript{(39)}、捲起塊\VUgras{窗簾} \textsuperscript{(41)}，將啲細行李擺上個
\VUgras{行李網} \textsuperscript{(43)}，然之後就坐低喺張\VUgras{座位} \textsuperscript{(42)} 度。
見到個\VUgras{點燈嘅} \textsuperscript{(44)} 嚟着燈，我先諗起自己要喺車上瞓一晚，
就叫住個負責租\VUgras{枕頭} \textsuperscript{(45)} 嘅職員，問佢攞一個。

%%K t12-11-1 p
11. --- 列車員嚟查飛。我問佢點解仲未開車，明明夠鐘㗎喇。佢答我話個
\VUgras{列車長} \textsuperscript{(46)} 仲喺度擺緊啲單車入\VUgras{行李卡}
\textsuperscript{(47)}。何況啲\VUgras{暖腳器} \textsuperscript{(48)} 都仲未換晒。

%%K t12-12-1 p
12. --- 不過就快開嘞：個\VUgras{司爐} \textsuperscript{(50)} 喺自己架\VUgras{煤水車}
\textsuperscript{(49)} 度剷煤，去谷旺\VUgras{火車頭} \textsuperscript{(54)} 嘅火；
\VUgras{個司機} \textsuperscript{(51)} 就淨係等\VUgras{副站長} \textsuperscript{(52)} 打訊號
——副站長就企喺\VUgras{月台} \textsuperscript{(53)} 上面。

%%K t12-13-1 p
13. --- 火車頭個\VUgras{鍋爐} \textsuperscript{(56)} 悶聲噉隆隆響；條\VUgras{煙囪}
\textsuperscript{(55)} 噴出一道道煙，撈埋啲\VUgras{蒸氣} \textsuperscript{(60)}。刺耳嘅
\VUgras{汽笛} \textsuperscript{(59)} 大大聲響，啲\VUgras{活塞} \textsuperscript{(58)} 動起嚟，
推住成部重機器嘅\VUgras{車輪} \textsuperscript{(57)}。

%%K t12-tit-3 sub
\VUsaut{3.0mm}
\VUpk{10.2pt}{40}{開車！}
\VUsaut{3.0mm}

%%K t12-14-1 p
14. --- 架車喺\VUgras{路軌} \textsuperscript{(64)} 上面越走越快；啲\VUgras{月台}
\textsuperscript{(65)} 唔見咗；我哋掠過啲\VUgras{廁所} \textsuperscript{(66)}，仲掠過

%%K t12-noto-1 noto
\VUnotes{143.2mm}{%
\textit{(*) 註。} --- \textit{呢段行程當然係照大戰之前嘅國界講嘅。}%
}

%%K t12-14-1 p suite
一列停咗喺另一條\VUgras{路軌} \textsuperscript{(63)} 上面嘅車：\VUgras{臥車}
\textsuperscript{(67)}、\VUgras{餐卡} \textsuperscript{(68)}、\VUgras{郵政卡}
\textsuperscript{(69)}。

%%K t12-15-1 p
15. --- 跟住，個\VUgras{貨站} \textsuperscript{(71)} 喺我哋眼前掠過。棚廠底下，
啲職員裝緊搭慢車走嘅\VUgras{貨} \textsuperscript{(72)}。有班\VUgras{調車工}
\textsuperscript{(76)} 喺個\VUgras{水塔} \textsuperscript{(73)} 附近調緊啲\VUgras{牲口卡}
\textsuperscript{(75)} 同啲\VUgras{平板卡} \textsuperscript{(79)}，砌一列\VUgras{貨車}
\textsuperscript{(80)} 出嚟。

%%K t12-16-1 p
16. --- 我哋過咗最後一個\VUgras{分岔} \textsuperscript{(78)}，扳道工就企喺隔籬；
我哋掠過個\VUgras{訊號牌} \textsuperscript{(86)}，個站就喺後面愈行愈遠，冇埋。

%%K t12-17-1 p
17. --- 好快我哋就過咗條\VUgras{鐵橋} \textsuperscript{(74)}，橋跨住條\VUgras{大河}
\textsuperscript{(81)}。過完橋，我哋沿住條河行；河面上有啲夠力嘅\VUgras{拖輪}
\textsuperscript{(82)} 拖住笨重嘅\VUgras{貨艇} \textsuperscript{(83)}。我哋仲見到一隻
\VUgras{客船} \textsuperscript{(84)}，佢啱啱埋緊個\VUgras{躉船碼頭} \textsuperscript{(85)}。

%%K t12-18-1 p
18. --- 飛快噉掠過幾個站之後，我哋穿過幾條\VUgras{隧道} \textsuperscript{(87)}，
拋低後面數唔清嘅\VUgras{屋仔} \textsuperscript{(88)}——係啲\VUgras{看柵佬}住嘅。
畀架車搖到似搖籃噉，我就瞓咗個大覺。

%%K t12-tit-4 sub
\VUsaut{3.0mm}
\VUpk{10.2pt}{40}{德意志－阿夫里庫爾站。}
\VUsaut{3.0mm}

%%K t12-19-1 p
19. --- 我畀一把職員嘅聲嗌醒：「德意志－阿夫里庫爾！\textsuperscript{(*)}
全部落車！」跟住就係海關檢查，同埋過境嗰堆煩到死嘅手續。我要照個站嘅
\VUgras{大鐘} \textsuperscript{(89)} 校返自己隻錶——個鐘快咗五十五分鐘；

%%K t12-19-1 p suite
啱啱瞓醒，我睇極都分唔清邊條係\VUgras{長針} \textsuperscript{(90)}、邊條係
\VUgras{短針} \textsuperscript{(91)}；連個\VUgras{錶面} \textsuperscript{(92)} 都畀朝早嗰陣霧
遮到矇晒。

%%K t12-20-1 p
20. --- 啲關員查嘢嗰陣，我一路打喊露一路讀住牆上面啲\VUgras{海報}
\textsuperscript{(93)}。我食咗個早餐消磨時間，又查咗德國嗰張\VUgras{鐵路網}
\textsuperscript{(94)} 圖，睇下仲差幾遠先到史特拉斯堡；然之後就返返上車卡，
心諗好快就到目的地，精神都好返晒。

\VUsaut{6.0mm}
\VUfilet{22mm}
